% Section "Related Work" (owner: related; audit framing). No result appears here.
% Cited keys exist in references.bib; each work is described only as far as CITATION_MAP.md supports.
% Cross-reference used from another section: lem:invariance (preliminaries.tex).
\section{Related Work}
\label{sec:related}

\paragraph{LLM-based recommendation and yes/no scoring.}
Pointwise LLM recommenders ask a binary question about a user--item pair and score the item by the probability of the
answer \emph{Yes}. \mbox{TALLRec}~\citep{bao2023tallrec} tunes an LLM with LoRA on such questions, CoLLM~\citep{zhang2025collm}
and \mbox{BinLLM}~\citep{zhang2024binllm} add collaborative information to the same formulation, and
ReLLa~\citep{lin2024rella} retrieves behaviours for CTR-style prediction. TALLRec, CoLLM and BinLLM evaluate \py{} by AUC
or its per-user average (UAUC), and we are not aware of an analysis of what the probability encodes. Other LLM
recommenders generate or select items~\citep{geng2022p5,liao2024llara,kim2024allmrec,bao2025bigrec};
D3~\citep{bao2024d3} shows that ghost tokens and length normalisation distort their sequence-level confidence, which a
single \texttt{Yes}/\texttt{No} answer position avoids, and S-DPO~\citep{chen2024sdpo} trains them with a softmax
preference loss over several negatives.
% METHOD-SLOT-BEGIN
Unlike CoLLM and BinLLM, the optional prior-offset LoRA adds a stop-gradient item-prior offset to the logit.
% METHOD-SLOT-END
We audit this pointwise family, zero-shot and LoRA-tuned, against labels that match the question and with non-LLM
reference rows in every table.

\paragraph{Confidence and calibration.}
Black-box elicitation finds LLMs overconfident when they verbalise confidence and separates calibration from failure
prediction~\citep{xiong2024llmuncertainty}; token probabilities of language models are fairly calibrated on
multiple-choice questions~\citep{kadavath2022language}; Conf\-Tuner trains verbalised confidence with a tokenised Brier score, a proper
scoring rule~\citep{li2025conftuner}; and proximity-informed calibration finds low-density samples more overconfident at
matched confidence~\citep{xiong2023procal}. Reliability is summarised by the expected calibration error
(ECE)~\citep{naeini2015ece}; temperature scaling~\citep{guo2017calibration} is strictly increasing and cannot reorder one
user's candidates (Lemma~\ref{lem:invariance}). For LLM recommenders, UQRec~\citep{kweon2025uqrec} decomposes list-wise
predictive uncertainty into recommendation and prompt uncertainty and adapts prompts; UGR~\citep{fan2026ugr} penalises
confident errors in training and uses confidence tokens for user rejection and item truncation in generative
recommendation; EviRank~\citep{yan2026evirank} estimates position-level confidence for LLM ranking. An audit of verbalised
confidence in zero-shot LLM recommenders is stratified by popularity and scored against catalogue
membership~\citep{ravikumar2026hallucinating}; KnowSA~\citep{lee2026knowsa} finds popularity and uncertainty poor
proxies of what an LLM knows about an item; a workshop study finds request-level selective escalation to an LLM
unhelpful in a sparse fixed-candidate setting~\citep{zhou2026routerec}. We instead audit token-level yes/no confidence of
pointwise recommenders under a registered protocol with matched labels, non-LLM references and fresh-user
confirmation, and test serving the most confident requests (selective classification~\citep{geifman2017selective}) against
random coverage.

\paragraph{Yes-bias, answer priors and item familiarity.}
Binary answers of language models carry input-independent biases: contextual calibration estimates the answer prior
from content-free inputs and removes it~\citep{zhao2021calibrate}, domain-conditional PMI counters surface-form
competition by dividing out an answer's prior~\citep{holtzman2021surface}, and batch calibration estimates the label
prior from a batch~\citep{zhou2024batch}. Yes/no studies report a global lean toward \emph{no} rather than human-like
acquiescence~\citep{braun2025acquiescence}, a framing asymmetry between positive and negative
formulations~\citep{zhang2025yesharder}, and an order- and wording-driven bias with idiosyncratic per-item
residuals~\citep{huang2026yesno}. Contrast-consistent search combines the probability of a statement with that of its
negation~\citep{burns2023ccs}; a two-view contrast of a question and its mirror belongs to this family and is not
claimed as new: it enters only as a registered arm (\method{}) that must beat a paraphrase placebo and a
correlation-matched ensemble null. LLMs also prefer recognised brands over fictional ones with identical
specifications~\citep{chu2026incumbent}. These works characterise a bias of the answer or of a brand; we ask how much of
a \emph{personalised} recommender's \py{} is a user-independent item prior rather than personal evidence, via a
reliability-corrected personal share, information gain over a prior-only item mean and a pseudonym knockout.

\paragraph{Popularity, calibrated recommendation and feedback loops.}
Popularity bias~\citep{abdollahpouri2017controlling} has causal treatments that subtract or deconfound the
item-popularity effect~\citep{wei2021macr,zhang2021pda}; in LLM recommenders, SFT and DPO amplify
it~\citep{gao2025flower,gao2025sprec}. Calibrated recommendation~\citep{steck2018calibrated} matches the genre mix of a
list to a user's history; the probability calibration used here instead follows work on calibrating ranking-model
scores under missing-not-at-random feedback~\citep{kweon2022calibrated}. Feedback loops raise
homogeneity~\citep{chaney2018feedback} and amplify popularity bias~\citep{mansoury2020feedback}, and LLM-agent
simulators support filter-bubble studies~\citep{zhang2024agent4rec}. Whereas these works concern the exposure produced by
a ranking score, we ask whether the confidence attached to a score is itself popularity-linked (its level and its
matched-confidence calibration gap, with popularity partialled on item quality), and we measure exposure statically on
same-candidate next-item panels, an existing benchmark protocol on which C-CRP, a verbalised-probability pointwise
reranker, and eight published LLM-enhanced recommenders serve as context rows only; no closed loop is simulated.

\paragraph{Exposure-biased evaluation.}
Logged feedback is missing not at random: propensity-based debiasing was introduced with the Coat test
set~\citep{schnabel2016treatments}, Yahoo!\ R3 supplies a randomly exposed test set~\citep{marlin2009collaborative}, and
KuaiRec a fully observed matrix~\citep{gao2022kuairec}. We avoid the missing-as-negative convention in calibration by
using explicit ratings whose threshold matches the question (rating $\ge4$ versus $\le2$). The zero-shot pilot
on KuaiRec missed its registered bar, so fully observed labels are not an endpoint; next-item panels, with sampled
negatives, are used only under same-candidate scope.

\paragraph{Data pruning for LLM-based recommendation.}
Data selection for LLM-recommender fine-tuning uses influence and effort scores~\citep{lin2024dealrec} or
optimal-transport coresets, where gradient- and loss-based scores fall below random selection on TALLRec-style CTR
tuning~\citep{mei2025goracs}; instruction tuning has cut claims the model is unfamiliar with at an uncertainty
quantile~\citep{wu2025unit}. We ask whether a recommender's own confidence identifies noisy training pairs, against
matched-random pruning over five seeds, and report the answer either way.
